\SourceLeaf{266}
\Needspace{8\baselineskip}
\Article{242}
The preceding propositions cover the composition of forms in its full generality. We now move on to more particular applications, with which we did not want to interrupt their sequence. First we return to the problem of Article 236, restricting it by the following conditions: \textit{first}, the forms to be composed have the same determinant, that is, $d=d'$; \textit{second}, $m$. $m'$ are coprime; \textit{third}, the required form is composed directly from each of $f$, $f'$. Then $mm$, $m'm'$ are also coprime, so the greatest common divisor of $dm'm'$, $d'mm$, \textit{that is,}\ $D$, becomes $=d=d'$, and $n=n'=1$. We set the four quantities $\mathfrak{Q}$, $\mathfrak{Q}'$, $\mathfrak{Q}''$, $\mathfrak{Q}'''$, which can be chosen at will, equal to $=-1$, $0$, $0$, $0$ respectively. This is always allowed except in the single case where $a$, $a'$, $b+b'$ are all $=0$, which we therefore leave aside here. Clearly this case can occur only for forms whose determinant is a positive square. It is then clear that $\mu$ becomes the greatest common divisor of $a$, $a'$, $b+b'$; the numbers $\mathfrak{P}'$, $\mathfrak{P}''$, $\mathfrak{P}'''$ must be chosen so that
\[
\mathfrak{P}'a+\mathfrak{P}''a'+\mathfrak{P}'''(b+b')=\mu
\]
while $\mathfrak{P}$ itself is entirely arbitrary. Substituting their values for $p$, $q$, $p'$, $q'$ etc. in the cited passage gives
\[
\begin{gathered}A=\frac{aa'}{\mu\mu},\\ B=\frac{1}{\mu}(\mathfrak{P}aa'+\mathfrak{P}'ab'+\mathfrak{P}''a'b+\mathfrak{P}'''(bb'+D))\end{gathered}
\]
and $C$ can be determined from the equation $AC=BB-D$, unless $a$ is zero and $a'=0$ also holds.

In this solution, then, the value of $A$ does not depend on the values of $\mathfrak{P}$, $\mathfrak{P}'$, $\mathfrak{P}''$, $\mathfrak{P}'''$ (which can be determined in infinitely many different ways). But $B$ takes different values when these numbers are assigned different values, and it is worth investigating how all the values of $B$ are connected. To this end we observe:

I. However $\mathfrak{P}$, $\mathfrak{P}'$, $\mathfrak{P}''$, $\mathfrak{P}'''$ are determined, the resulting values of $B$ are all congruent modulo $A$. Suppose that setting
\[
\begin{gathered}\mathfrak{P}=\mathfrak{p},\\ \mathfrak{P}'=\mathfrak{p}',\\ \mathfrak{P}''=\mathfrak{p}'',\\ \mathfrak{P}'''=\mathfrak{p}''',\\\text{gives } B=\mathfrak{B}\end{gathered}
\]
but setting
\[
\begin{gathered}\mathfrak{P}=\mathfrak{p}+\mathfrak{d},\\ \mathfrak{P}'=\mathfrak{p}'+\mathfrak{d}',\\ \mathfrak{P}''=\mathfrak{p}''+\mathfrak{d}''.\\ \mathfrak{P}'''=\mathfrak{p}'''+\mathfrak{d}''',\\\text{gives } B=\mathfrak{B}+\mathfrak{D}\end{gathered}
\]
\SourceLeaf{267}
Then
\[
\begin{gathered}a\mathfrak{d}'+a'\mathfrak{d}''+(b+b')\mathfrak{d}'''=0,\\
aa'\mathfrak{d}+ab'\mathfrak{d}'+a'b\mathfrak{d}''+(bb'+D)\mathfrak{d}'''=\mu\mathfrak{D}\end{gathered}
\]
Multiply the first side of the latter equation by $a\mathfrak{p}'+a'\mathfrak{p}''+(b+b')\mathfrak{p}'''$, the second by $\mu$, and subtract from the first product the quantity
\[
(ab'\mathfrak{p}'+a'b\mathfrak{p}''+(bb'+D)\mathfrak{p}''')(a\mathfrak{d}'+a'\mathfrak{d}''+(b+b')\mathfrak{d}''')
\]
which, by the former equation, is clearly $=0$. Expanding and cancelling the terms that cancel each other gives
\[
\begin{aligned}aa'(&\mu\mathfrak{d}\\&+((b'-b)\mathfrak{p}''+c'\mathfrak{p}''')\mathfrak{d}'\\&+((b-b')\mathfrak{p}'+c\mathfrak{p}''')\mathfrak{d}''\\&-(c'\mathfrak{p}'+c\mathfrak{p}'')\mathfrak{d}''')=\mu\mu\mathfrak{D}\end{aligned}
\]
so clearly $\mu\mu\mathfrak{D}$ is divisible by $aa'$, or $\mathfrak{D}$ by $\frac{aa'}{\mu\mu}$, \textit{that is,}\ by $A$, and
\[
\mathfrak{B}\equiv\mathfrak{B}+\mathfrak{D}\ (\mathrm{mod.}\ A)
\]

II. If the values $\mathfrak{p}$, $\mathfrak{p}'$, $\mathfrak{p}''$ $\mathfrak{p}'''$ of $\mathfrak{P}$, $\mathfrak{P}'$, $\mathfrak{P}''$, $\mathfrak{P}'''$ give $B=\mathfrak{B}$, other values of these numbers can be found that give $B$ any specified value congruent to $\mathfrak{B}$ modulo $A$, say $\mathfrak{B}+kA$. First we observe that the four numbers $\mu$, $c$, $c'$, $b-b'$ cannot have a common divisor. For if they had one, it would divide the six numbers $a$, $a'$, $b+b'$, $c$, $c'$, $b-b'$, and therefore both $a$, $2b$, $c$ and $a'$, $2b'$, $c'$, and consequently also $m$, $m'$, which are coprime by hypothesis. We can therefore assign four integers $h$, $h'$, $h''$, $h'''$ so that
\[
h\mu+h'c+h''c'+h'''(b-b')=1
\]
Having done this, if we set
\[
\begin{aligned}
kh&=\mathfrak{d},\\k(h''(b+b')-h'''a')&=\mu\mathfrak{d}'\\
k(h'(b+b')+h'''a)&=\mu\mathfrak{d}'',\\-k(h'a'+h''a)&=\mu\mathfrak{d}'''
\end{aligned}
\]
it is clear that $\mathfrak{d}$, $\mathfrak{d}'$, $\mathfrak{d}''$, $\mathfrak{d}'''$ are integers. Further, we easily confirm that
\[
a\mathfrak{d}'+a'\mathfrak{d}''+(b+b')\mathfrak{d}'''=0
\]
\[
\begin{aligned}aa'\mathfrak{d}+ab'\mathfrak{d}'+a'b\mathfrak{d}''+(bb'+D)\mathfrak{d}'''\\&=\frac{aa'k}{\mu}(\mu h+ch'+c'h''+(b-b')h''')\\&=\mu kA\end{aligned}
\]
The former equation shows that $\mathfrak{p}+\mathfrak{d}$, $\mathfrak{p}'+\mathfrak{d}'$, $\mathfrak{p}''+\mathfrak{d}''$, $\mathfrak{p}'''+\mathfrak{d}'''$ are also values of $\mathfrak{P}$, $\mathfrak{P}'$, $\mathfrak{P}''$, $\mathfrak{P}'''$; the latter shows that these values produce $B=\mathfrak{B}+kA$. \textit{Q. E. D.} ---
\SourceLeaf{268}
It is clear from this that $B$ can always be determined so that it lies between $0$ and $A-1$ inclusive if $A$ is positive, or between $0$ and $-A-1$ if $A$ is negative.

\Needspace{8\baselineskip}
\Article{243}
From the equations
\[
\begin{gathered}\mathfrak{P}'a+\mathfrak{P}''a'+\mathfrak{P}'''(b+b')=\mu,\\
B=\frac{1}{\mu}(\mathfrak{P}aa'+\mathfrak{P}'ab'+\mathfrak{P}''a'b+\mathfrak{P}'''(bb'+D))\end{gathered}
\]
we obtain
\[
\begin{aligned}B&=b+\frac{a}{\mu}(\mathfrak{P}a'+\mathfrak{P}'(b'-b)-\mathfrak{P}'''c)\\&=b'+\frac{a'}{\mu}(\mathfrak{P}a+\mathfrak{P}''(b-b')-\mathfrak{P}'''c')\end{aligned}
\]
so
\[
B\equiv b\ (\mathrm{mod.}\ \tfrac{a}{\mu})\quad\text{and}\quad B\equiv b'\ (\mathrm{mod.}\ \tfrac{a'}{\mu})
\]
Whenever $\frac{a}{\mu}$, $\frac{a'}{\mu}$ are coprime, there is only one number between $0$ and $A-1$ (or between $0$ and $-A-1$ when $A$ is negative) that modulo $\frac{a}{\mu}$ is $\equiv b$, and is $\equiv b'$ modulo $\frac{a'}{\mu}$. If we set this number $=B$ and $\frac{BB-D}{A}=C$, it is clear that $(A,B,C)$ is composed from the forms $(a,b,c)$, $(a',b',c')$. In this case, then, to find the composite form we no longer need to consider the numbers $\mathfrak{P}$, $\mathfrak{P}'$, $\mathfrak{P}''$, $\mathfrak{P}'''$\HistoricalNote{*)}{This can always be done by using the congruences
\[
\frac{aB}{\mu}\equiv\frac{ab'}{\mu},\quad\frac{a'B}{\mu}\equiv\frac{a'b}{\mu},\quad\frac{(b+b')B}{\mu}\equiv\frac{bb'+D}{\mu}\ (\mathrm{mod.}\ A)
\]
}. Thus, \textit{for example,}\ if we seek a form composed from $(10,3,11)$, $(15,2,7)$, then $a$, $a'$, $b+b'$ are respectively $=10,15,5$; $\mu=5$; so $A=6$; $B\equiv3\ (\mathrm{mod.}\ 2)$ and $\equiv2\ (\mathrm{mod.}\ 3)$, giving $B=5$ and $(6,5,21)$ as the required form. --- The condition that $\frac{a}{\mu}$, $\frac{a'}{\mu}$ are coprime is entirely equivalent to the condition that the two numbers $a$, $a'$ have no greater common divisor than the three $a$, $a'$, $b+b'$; or, what comes to the same thing, that the greatest common divisor of $a$, $a'$ also divides $b+b'$. The following particular cases should especially be noted:

1) Given two forms $(a,b,c)$, $(a',b',c')$ of the same determinant $D$, such that the greatest common divisor of $a$, $2b$, $c$ is coprime to the greatest common divisor of $a'$, $2b'$, $c'$, and $a$ is coprime to $a'$, the form $(A,B,C)$ composed from them is found by setting $A=aa'$, $B\equiv b\ (\mathrm{mod.}\ a)$ and $\equiv b'\ (\mathrm{mod.}\ a')$, $C=\frac{BB-D}{A}$. This case always applies when one of the forms to be composed is the principal form, say $a=1$, $b=0$, $c=-D$. Then $A=a'$, and $B$ can be set
\SourceLeaf{269}
$=b'$, giving $C=c'$; so \textit{the form composed from the principal form and any other form of the same determinant is that other form itself.}

2) If two \textit{opposite}, properly primitive forms, say $(a,b,c)$ and $(a,-b,c)$, are to be composed, then $\mu=a$. It is easy to see from this that the principal form $(1,0,-D)$ is composed from them.

3) Given any number of properly primitive forms $(a,b,c)$, $(a',b',c')$, $(a'',b'',c'')$ etc. of the same determinant $D$, whose leading coefficients $a$, $a'$, $a''$ etc. are pairwise coprime, the form $(A,B,C)$ composed from all of them is found. Set $A$ equal to the product of all $a$, $a'$, $a''$ etc.; $B$ congruent to $b$, $b'$, $b''$ etc. modulo $a$, $a'$, $a''$ etc. respectively; $C=\frac{BB-D}{A}$. For it is easy to see that from the two forms $(a,b,c)$, $(a',b',c')$ the form $(aa',B,\frac{BB-D}{aa'})$ is composed; from this and $(a'',b'',c'')$ the form $(aa'a'',B,\frac{BB-D}{aa'a''})$ etc. --- Conversely:

4) Given a properly primitive form $(A,B,C)$ of determinant $D$, resolve the coefficient $A$ into any number of pairwise coprime factors $a$, $a'$, $a''$ etc.; choose numbers $b$, $b'$, $b''$ etc. equal to $B$, or at least congruent to it modulo $a$, $a'$, $a''$ etc. respectively, and set $ac=bb-D$, $a'c'=b'b'-D$, $a''c''=b''b''-D$ etc. Then the form $(A,B,C)$ is composed from the forms $(a,b,c)$, $(a'\ b'\ c')$, $(a''\ b''\ c'')$, or is \textit{resolvable into these forms}. It is proved without difficulty that the same proposition still holds even if the form $(A,B,C)$ is improperly primitive or derived. In this way, then, any form can be resolved into others of the same determinant whose leading coefficients are all either primes or powers of primes. Such a resolution can often be usefully applied when one form is to be composed from several given forms. Thus, \textit{for example,}\ if we seek the form composed from $(3,1,134)$, $(10,3,41)$, $(15,2,27)$, resolve the second into $(2,1,201)$, $(5,-2,81)$, and the third into $(3,-1,134)$, $(5,2,81)$. Clearly the form composed from the five forms $(3,1,134)$, $(2,1,201)$, $(5,-2,81)$, $(3,-1,134)$, $(5,2,81)$, in whatever order they are taken, is also composed from the three given forms. But composing the first with the fourth gives the principal form $(1,0,401)$; composing the third with the fifth gives the same form; so composing all of them gives the form $(2,1,201)$.

\SourceLeaf{270}
5) Since the method is useful, it is worth explaining it further. The preceding observation shows that the problem of composing any number of given properly primitive forms of the same determinant can be reduced to composing forms whose leading coefficients are powers of primes (for a prime can be regarded as the first power of itself). We should therefore consider especially the case where two properly primitive forms $(a,b,c)$, $(a',b',c')$ are to be composed, in which $a$ and $a'$ are powers of the \textit{same} prime. Let $a=h^{\varkappa}$, $a'=h^{\lambda}$, where $h$ denotes a prime, and suppose, as we may, that $\varkappa$ is not less than $\lambda$. Then $h^{\lambda}$ is the greatest common divisor of $a$, $a'$. If it also divides $b+b'$, we have the case considered at the beginning of this article, and $(A,B,C)$ is composed from the given forms if we set $A=h^{\varkappa-\lambda}$, $B\equiv b\ (\mathrm{mod.}\ h^{\varkappa-\lambda})$ and $\equiv b'\ (\mathrm{mod.}\ 1)$, where the latter condition can clearly be omitted; $C=\frac{BB-D}{A}$. But if $h^{\lambda}$ does not divide $b+b'$, the greatest common divisor of these numbers must itself be a power of $h$. Let it be $=h^{\nu}$; then $\nu<\lambda$ (we must set $\nu=0$. if $h^{\lambda}$ and $b+b'$ happen to be coprime). If, then, $\mathfrak{P}'$, $\mathfrak{P}''$, $\mathfrak{P}'''$ are determined so that
\[
\mathfrak{P}'h^{\varkappa}+\mathfrak{P}''h^{\lambda}+\mathfrak{P}'''(b+b')=h^{\nu}
\]
while $\mathfrak{P}$ is chosen at will, the form $(A,B,C)$ is composed from the given forms if we set
\[
\begin{gathered}A=h^{\varkappa+\lambda-2\nu},\\
B=b+h^{\varkappa-\nu}(\mathfrak{P}h^{\lambda}-\mathfrak{P}'(b-b')-\mathfrak{P}'''c),\\
C=\frac{BB-D}{A}\end{gathered}
\]
But it is easy to see that in this case $\mathfrak{P}'$ can also be chosen at will. So setting $\mathfrak{P}=\mathfrak{P}'=0$ gives
\[
B=b-\mathfrak{P}'''ch^{\varkappa-\nu}
\]
or, in general,
\[
B=kA+b-\mathfrak{P}'''ch^{\varkappa-\nu}
\]
where $k$ denotes an arbitrary number (preceding article). Only $\mathfrak{P}'''$ enters this very simple formula; it is a value of the expression $\frac{h^{\nu}}{b+b'}\ (\mathrm{mod.}\ h^{\lambda})$\HistoricalNote{*)}{or of the expression $\frac{1}{\frac{b+b'}{h^{\nu}}}\ (\mathrm{mod.}\ h^{\lambda-\nu})$, giving $B\equiv b-\frac{ch^{\varkappa-\nu}}{\frac{b+b'}{h^{\nu}}}\equiv\frac{(D+bb'):h^{\nu}}{(b+b'):h^{\nu}}\ (\mathrm{mod.}\ A)$}. If, \textit{for example,}\ we seek the form composed from $(16,3,19)$ and $(8,1,37)$, then $h=2$, $\varkappa=4$, $\lambda=3$, $\nu=2$.
\SourceLeaf{271}
So $A=8$, and $\mathfrak{P}'''$ is a value of the expression $\frac{4}{4}\ (\mathrm{mod.}\ 8)$, such as $1$, giving $B=8k-73$. Thus, setting $k=9$ gives $B=-1$ and $C=37$, or $(8,-1,37)$ as the required form.

Given, then, any number of forms whose leading coefficients are all powers of primes, we must check whether some of their leading coefficients are powers of the \textit{same} prime and compose these with each other using the rule just given. This produces forms whose leading coefficients are still powers of primes, but of entirely different primes; the form composed from them can therefore be determined by the third observation. \textit{For example,}\ given the forms $(3,1,47)$, $(4,0,35)$, $(5,0,28)$, $(16,2,9)$, $(9,7,21)$, $(16,6,11)$, the first and fifth give the form $(27,7,7)$; the second and fourth give $(16,-6,11)$, and this and the sixth give $(1,0,140)$, which can be left out. This leaves $(5,0,28)$, $(27,7,7)$, which produce $(135,-20,4)$; in its place we can take the properly equivalent $(4,0,35)$. This, then, is the result of composing the six given forms.

Other useful practical methods can also be drawn from this source. But to avoid going on too long, we leave out a fuller treatment of this matter and move on to more difficult things.
